\chapter{Recurrent Networks \& Sequence Modeling}\label{ch:rnn}
\begin{center}\small\textit{Language, speech, and time series unfold step by step -- networks must remember.}\end{center}

\noindent\fbox{\parbox{\dimexpr\linewidth-2\fboxsep-2\fboxrule}{%
\textbf{From zero: no sequence background assumed.} We start from why fixed windows fail on language and build from Elman RNNs to LSTM/GRU gating. Every formula is preceded by a picture; hidden states are defined on first use.}}
\vspace{0.3em}
\section*{Learning Outcomes}
After this chapter you will be able to:
\begin{enumerate}[label=(L\arabic*)]
    \item write the Elman RNN recurrence and its unrolling through time;
    \item explain vanishing and exploding gradients in vanilla RNNs;
    \item write LSTM/GRU gate equations and explain what each gate does;
    \item compare bidirectional/stacked RNNs and the bridge to attention.
\end{enumerate}

\section{Intuition: Memory Through Loops}
An MLP sees a fixed window; language does not fit. ``The cat that chased the dogs \emph{was}...'' -- agreement depends on ``cat'' six words back. We need a network with \textbf{memory}: a hidden state that carries information forward. A \textbf{Recurrent Neural Network (RNN)} does exactly this -- it reuses the same weights at each time step, updating a state $\mathbf{h}_t$ from input $\mathbf{x}_t$ and previous state $\mathbf{h}_{t-1}$.

Unrolling the loop in time reveals a deep network whose depth is the sequence length -- with shared weights and a gradient that must flow back through time.

\section{The Vanilla RNN}
\begin{definition}[Elman RNN]
For sequence $\mathbf{x}_{1:T}$,
\begin{equation}
\mathbf{h}_t = \tanh(W_{hh}\mathbf{h}_{t-1}+W_{xh}\mathbf{x}_t+\mathbf{b}_h),\quad \mathbf{y}_t = W_{hy}\mathbf{h}_t+\mathbf{b}_y,\quad \mathbf{h}_0=\mathbf{0}
\end{equation}
Loss sums over time: $\mathcal{L}=\sum_t \ell(\mathbf{y}_t,\mathbf{y}^*_t)$.
\end{definition}

\begin{center}
\begin{tikzpicture}[font=\scriptsize, >=Stealth, scale=0.93]
  % Unrolled RNN
  \foreach \t in {1,2,3} {
    \node[draw, fill=blue!12, rounded corners, minimum width=1.1cm, minimum height=0.7cm] (x\t) at (2*\t, -1.2) {$\mathbf{x}_{\t}$};
    \node[draw, fill=green!15, circle, minimum size=1.0cm, thick] (h\t) at (2*\t, 0.4) {$\mathbf{h}_{\t}$};
    \node[draw, fill=orange!12, rounded corners, minimum width=1.0cm, minimum height=0.6cm] (y\t) at (2*\t, 1.8) {$\mathbf{y}_{\t}$};
    \draw[->, thick, blue!60!black] (x\t) -- (h\t);
    \draw[->, thick, orange!70!black] (h\t) -- (y\t);
  }
  \draw[->, thick, green!50!black, bend left=12] (h1) -- (h2);
  \draw[->, thick, green!50!black, bend left=12] (h2) -- (h3);
  \node[left, font=\tiny, text=green!50!black] at (3,0.9) {$W_{hh}$ shared};
  \draw[->, thick, red!60!black, dashed] (y3) -- ++(0.8,0) |- (h1) node[midway, above, font=\tiny, text=red!60!black] {BPTT};
  \node[draw, fill=violet!8, rounded corners, font=\tiny, align=center] at (4,-2.3) {Unrolled over time\\depth = sequence length};
\end{tikzpicture}
\end{center}

\textbf{Backpropagation Through Time (BPTT)}: unroll, then backprop. Gradient involves repeated multiplication by $W_{hh}$:
\begin{equation}
\frac{\partial\mathcal{L}_T}{\partial\mathbf{h}_1}=\frac{\partial\mathcal{L}_T}{\partial\mathbf{h}_T}\prod_{t=2}^{T} W_{hh}^{\top}\operatorname{diag}(\tanh'(\cdot))
\end{equation}
If spectral radius $<1$, gradient vanishes exponentially; if $>1$, it explodes. This is why vanilla RNNs cannot learn dependencies beyond $\sim$10 steps.

\section{LSTM and GRU: Gated Memory}
Gates -- learnable sigmoidal switches in $[0,1]$ -- protect the gradient by creating an additive ``conveyor belt''.

\begin{definition}[LSTM, Hochreiter \& Schmidhuber 1997]
\begin{align}
\mathbf{f}_t&=\sigma(W_f[\mathbf{h}_{t-1},\mathbf{x}_t]+\mathbf{b}_f) &&\text{forget gate}\\
\mathbf{i}_t&=\sigma(W_i[\mathbf{h}_{t-1},\mathbf{x}_t]+\mathbf{b}_i) &&\text{input gate}\\
\tilde{\mathbf{c}}_t&=\tanh(W_c[\mathbf{h}_{t-1},\mathbf{x}_t]+\mathbf{b}_c) &&\text{candidate}\\
\mathbf{c}_t&=\mathbf{f}_t\odot\mathbf{c}_{t-1}+\mathbf{i}_t\odot\tilde{\mathbf{c}}_t &&\text{cell state (additive!)}\\
\mathbf{o}_t&=\sigma(W_o[\mathbf{h}_{t-1},\mathbf{x}_t]+\mathbf{b}_o) &&\text{output gate}\\
\mathbf{h}_t&=\mathbf{o}_t\odot\tanh(\mathbf{c}_t)
\end{align}
Key: $\mathbf{c}_t$ updates \emph{additively}; gradient $\partial\mathbf{c}_T/\partial\mathbf{c}_t=\prod \mathbf{f}_k$ stays near $1$ if forget gates stay open.
\end{definition}

\begin{definition}[GRU, Cho et al. 2014]
Simpler, two gates:
\begin{align}
\mathbf{z}_t&=\sigma(W_z[\mathbf{h}_{t-1},\mathbf{x}_t]),\quad \mathbf{r}_t=\sigma(W_r[\mathbf{h}_{t-1},\mathbf{x}_t])\\
\tilde{\mathbf{h}}_t&=\tanh(W_h[\mathbf{r}_t\odot\mathbf{h}_{t-1},\mathbf{x}_t]),\quad \mathbf{h}_t=(1-\mathbf{z}_t)\odot\mathbf{h}_{t-1}+\mathbf{z}_t\odot\tilde{\mathbf{h}}_t
\end{align}
Often matches LSTM with fewer parameters.
\end{definition}

\begin{center}
\begin{tikzpicture}[font=\tiny, >=Stealth, scale=0.88]
  % LSTM cell schematic
  \node[draw, fill=blue!10, rounded corners, minimum width=5.5cm, minimum height=3.0cm] (box) at (0,0) {};
  \node[above, font=\scriptsize, text=blue!60!black] at (0,1.5) {LSTM Cell};
  \draw[->, thick, green!50!black] (-3.2,0.8) -- (-1.6,0.8) node[midway, above] {$\mathbf{c}_{t-1}$};
  \draw[->, thick, green!50!black] (1.6,0.8) -- (3.2,0.8) node[midway, above] {$\mathbf{c}_t$};
  \node[draw, fill=orange!18, circle, minimum size=0.6cm] (forget) at (-0.8,0.3) {$\times$};
  \node[draw, fill=orange!18, circle, minimum size=0.6cm] (add) at (0.5,0.3) {$+$};
  \node[draw, fill=violet!15, rounded corners, minimum width=0.9cm, minimum height=0.5cm] (fg) at (-0.8,-0.7) {$\sigma$ forget};
  \node[draw, fill=violet!15, rounded corners, minimum width=0.9cm, minimum height=0.5cm] (ig) at (0.5,-0.7) {$\sigma$ input};
  \draw[->, thin, violet!60!black] (fg) -- (forget); \draw[->, thin, violet!60!black] (ig) -- (add);
  \node[below, font=\tiny, text=gray!60!black, align=center] at (0,-1.7) {Cell = conveyor belt\\gates control read/write};
  \node[draw, fill=red!12, circle, minimum size=0.5cm] (out) at (2.0,-0.3) {$\times$};
  \draw[->, thin, green!50!black] (-1.6,0.8) -- (forget) -- (add) -- (1.6,0.8);
  \draw[->, thin, blue!60!black] (add) -- (out);
\end{tikzpicture}
\end{center}

\section{Attention Intuition (Bridge to Transformers)}
RNNs compress all history into $\mathbf{h}_t$ -- a bottleneck. \textbf{Attention} lets the decoder peek at all encoder states: for each output step, compute weights $\alpha_{t,i}=\text{softmax}(\text{score}(\mathbf{s}_{t-1},\mathbf{h}_i))$ and context $\mathbf{c}_t=\sum_i\alpha_{t,i}\mathbf{h}_i$. The network \emph{learns where to look}. This idea, scaled up and stripped of recurrence, becomes the Transformer.

\section{Worked Example}
\begin{example}[Character-Level RNN in PyTorch]
\begin{lstlisting}[language=Python]
import torch, torch.nn as nn
class CharRNN(nn.Module):
    def __init__(self, vocab=65, hidden=128):
        super().__init__()
        self.emb = nn.Embedding(vocab, 64)
        self.rnn = nn.LSTM(64, hidden, batch_first=True)  # or nn.GRU
        self.fc  = nn.Linear(hidden, vocab)
    def forward(self, x, h=None):
        e = self.emb(x)          # (B,T,64)
        out, h = self.rnn(e, h)  # out: (B,T,hidden)
        return self.fc(out), h   # logits per step

# Training: teacher forcing
model = CharRNN(); opt = torch.optim.Adam(model.parameters(), lr=1e-3)
for x,y in loader:  # x:(B,T) y:(B,T) next-char targets
    opt.zero_grad()
    logits,_ = model(x)
    loss = nn.CrossEntropyLoss()(logits.view(-1,65), y.view(-1))
    loss.backward()
    nn.utils.clip_grad_norm_(model.parameters(), 1.0)  # cure explosion
    opt.step()

# Sampling
model.eval()
x = torch.tensor([[ord('T')%65]])
with torch.no_grad():
    for _ in range(50):
        logits,h = model(x[:,-1:])
        nxt = torch.multinomial(torch.softmax(logits[:,-1],-1),1)
        x = torch.cat([x,nxt],1)
print(''.join(chr(c.item()) for c in x[0]))
\end{lstlisting}
\end{example}

Bidirectional RNNs run forward and backward, concatenating states for tasks like NER where future context matters. Stacking 2--4 layers deepens representation.

\section{Why Vanilla RNNs Fail: Vanishing/Exploding Gradients}
Consider $\mathbf{h}_t = \tanh(W_{hh}\mathbf{h}_{t-1}+\dots)$. Then
$\partial\mathbf{h}_T/\partial\mathbf{h}_t = \prod_{k=t+1}^{T} D_k W_{hh}$ where
$D_k=\operatorname{diag}(\tanh'(\cdot))$ has entries in $(0,1]$. Norm behaves as
$\|W_{hh}\|^{T-t}$. If $\|W_{hh}\|<1$, product $\to0$ exponentially; if $>1$,
$\to\infty$. Gradient clipping handles explosion, but vanishing requires
architecture -- gates. Orthogonal init of $W_{hh}$ (eigenvalues $\approx1$) helps
but is insufficient alone for long horizons.

\section{LSTM Gates in Action}
Walk through a memory write/read: to remember a fact at $t=5$ for use at $t=50$,
the forget gate must stay $\approx1$ for 45 steps ($\mathbf{b}_f$ initialized to
$1$ helps), input gate opens at $t=5$ to write $\tilde{\mathbf{c}}_5$, cell
holds $\mathbf{c}_t$ additively, output gate opens at $t=50$ to expose it.
Peephole connections (optional) let gates see $\mathbf{c}$. GRU merges cell and
hidden, using update gate $\mathbf{z}_t$ to interpolate old vs new -- fewer
parameters, similar performance, faster.

\section{Bidirectional and Stacked RNNs}
A \textbf{Bi-RNN} runs forward $\overrightarrow{\mathbf{h}}_t$ and backward
$\overleftarrow{\mathbf{h}}_t$, concatenating $[\overrightarrow{\mathbf{h}}_t;
\overleftarrow{\mathbf{h}}_t]$ for tasks where future context helps (NER, speech).
\textbf{Stacked} (deep) RNNs feed $\mathbf{h}^{(l)}_t$ as input to layer $l+1$,
learning hierarchical timescales. Attention over encoder states (Bahdanau,
Luong) then lets decoders focus: $\alpha_{t,i}\propto\exp(\mathbf{s}_{t-1}^{\top}W\mathbf{h}_i)$,
$\mathbf{c}_t=\sum_i\alpha_{t,i}\mathbf{h}_i$, finally $\tilde{\mathbf{s}}_t=
\tanh(W_c[\mathbf{c}_t;\mathbf{s}_t])$ before predicting $y_t$.

\section{RNNs Today and the Path to Transformers}
LSTMs dominated 2014--2017 for translation, speech, and language modeling.
But sequential dependence prevents parallelization ($O(T)$ steps) and path length
between positions is $O(T)$. Attention shortens it to $O(1)$ and parallelizes
to $O(1)$ steps with $O(T^2)$ work -- the Transformer keeps the gating insight
(additive paths, normalization) but removes recurrence entirely.

\section{Generative Models: Autoencoders, VAEs, GANs, Diffusion}
\label{sec:generative}
Our CharRNN already \emph{generates}: given history it predicts the next character, and sampling turns those predictions into new text. Modern generative models generalize this idea -- learn the data distribution $p(\mathbf{x})$ itself, then sample brand-new $\mathbf{x}$ from it. Three landmarks cover the design space.
\paragraph{Autoencoder to VAE.} An \textbf{autoencoder} compresses $\mathbf{x}$ through a bottleneck $\mathbf{z}=f_{\mathrm{enc}}(\mathbf{x})$ and reconstructs $\hat{\mathbf{x}}=f_{\mathrm{dec}}(\mathbf{z})$, trained to minimize $\|\mathbf{x}-\hat{\mathbf{x}}\|^2$. A \textbf{variational autoencoder} (Kingma--Welling, 2013) makes $\mathbf{z}$ a distribution $q(\mathbf{z}\mid\mathbf{x})$ with prior $p(\mathbf{z})=\mathcal{N}(0,I)$, so we can sample $\mathbf{z}\sim p(\mathbf{z})$ and decode novel $\mathbf{x}$. Training maximizes the ELBO one-liner:
\begin{equation}
\log p(\mathbf{x})\geq \mathbb{E}_{q}[\log p(\mathbf{x}\mid\mathbf{z})]-\mathrm{KL}(q(\mathbf{z}\mid\mathbf{x})\,\|\,p(\mathbf{z})).
\label{eq:elbo}
\end{equation}
The first term pulls reconstructions toward realism; the KL term keeps the latent space dense and samplable -- without it, gaps in $\mathbf{z}$-space decode to garbage.
\paragraph{GANs.} A \textbf{GAN} (Goodfellow et al., 2014) pits a generator $G(\mathbf{z})$ against a discriminator $D(\mathbf{x})$ in a minimax game: $\min_G\max_D\,\mathbb{E}_{\mathbf{x}}[\log D(\mathbf{x})]+\mathbb{E}_{\mathbf{z}}[\log(1-D(G(\mathbf{z})))]$. $D$ learns to spot fakes, $G$ learns to fool $D$; at equilibrium $G$ draws from $p_{\mathrm{data}}$. Samples are sharp, but training can collapse to a few modes and offers no likelihood.
\paragraph{Diffusion.} A \textbf{diffusion model} (Ho et al., 2020) destroys structure by adding Gaussian noise over $T$ steps, then learns the reverse: a denoiser $\epsilon_\theta(\mathbf{x}_t,t)$ predicting the noise present at each step. Sampling starts from pure noise $\mathbf{x}_T\sim\mathcal{N}(0,I)$ and iteratively denoises down to $\mathbf{x}_0$. Slow (many steps) but stable, with the best image quality to date. Our CharRNN sampler is its autoregressive cousin -- both turn a seed (or noise) into data one refinement at a time.

\section{Common Pitfalls}
Feeding the whole sequence without truncation causes $O(T)$ memory blowup and
slow BPTT -- use truncated BPTT (backprop every $k$ steps, carry state). Not
detaching hidden states between batches leaks gradient across unrelated
sequences. Teacher forcing during training vs autoregressive sampling at test
causes exposure bias -- scheduled sampling or professor forcing mitigates.
Always clip gradients; without clipping, one bad batch can NaN the run.

\section{Further Reading}
Hochreiter \& Schmidhuber (1997) LSTM; Cho et al. (2014) GRU; Bahdanau et al.
(2015) attention. Graves (2013) applies RNNs to speech; Karpathy's ``The
Unreasonable Effectiveness of RNNs'' is the classic blog. For BPTT theory, see
Pascanu et al. (2013) on vanishing/exploding.

\section{Chapter Summary}
RNNs share weights over time: $\mathbf{h}_t=\tanh(W_{hh}\mathbf{h}_{t-1}+
W_{xh}\mathbf{x}_t)$. BPTT unrolls and backprops through the product of
Jacobians -- vanishing/exploding follows. LSTM/GRU gates create additive
memory ($\mathbf{c}_t=\mathbf{f}_t\odot\mathbf{c}_{t-1}+\mathbf{i}_t\odot
\tilde{\mathbf{c}}_t$) preserving gradient. Attention fixes the bottleneck by
peeking at all states, foreshadowing Transformers.

\smallskip
\noindent\textit{Key takeaways:} (1) recurrence $=$ memory via loop;
(2) BPTT $=$ backprop through time; (3) gates $=$ learnable switches;
(4) $\mathbf{f}\approx1$ preserves long memory; (5) attention $>$ compression.

\paragraph{Gradient sanity check.} Before training, run gradient checking with $\epsilon=10^{-5}$ on a 3-step unroll. If relative error $>10^{-3}$, check gate biases (forget bias $b_f=1$ helps) and clipping threshold.

\begin{example}[BPTT product] For $h_t=\tanh(W h_{t-1}+U x_t)$, $\partial h_3/\partial h_1 = \prod_{k=2}^3 \text{diag}(1-\tanh^2(\cdot))W$. Spectral radius $\rho(W)>1$ explodes, $\rho(W)<1$ vanishes. Gating replaces product with sum via $c_t$.
\end{example}

\section{Exercises}
\begin{exercise} Unroll a vanilla RNN for $T=3$ and derive $\partial\mathcal{L}_3/\partial W_{hh}$ explicitly. Identify the product that vanishes/explodes.\end{exercise}
\begin{exercise} Explain in words how LSTM's forget gate $\mathbf{f}_t\approx1$ preserves gradient over 100 steps. What initialization of $\mathbf{b}_f$ encourages this?\end{exercise}
\begin{exercise} Implement a NumPy forward pass for one LSTM step. Count parameters for $d_x=64$, $d_h=128$ and compare to GRU and vanilla RNN.\end{exercise}
\begin{exercise} Train CharRNN on Tiny Shakespeare (10k chars) with RNN vs LSTM (same hidden size). Plot loss and report which captures longer dependencies -- test by checking if it closes a quote 30 chars later.\end{exercise}
\begin{exercise} For seq2seq translation, why does attention help with long sentences? Sketch the alignment matrix for ``The cat sat'' $\to$ ``Le chat s'est assis''.\end{exercise}
% End of Chapter
